\chapter{Methodology}
\label{ch:method}

This chapter describes the construction we use to enforce set-level
flatness against a learned attacker. \Cref{sec:method-overview} states
the pipeline and its main design choices. \Cref{sec:method-formalism}
formalizes the attacker-induced distribution and the flatness metrics
computed on top of it. \Cref{sec:method-corpora} describes the corpora
and the disjoint oracle/attacker split. \Cref{sec:method-cnn} specifies
the character-level CNN that plays the dual role of defensive oracle
during filtering and adversary during evaluation. \Cref{sec:method-gens}
describes the ten honeyword generators we benchmark.
\Cref{sec:method-filter} presents the hybrid flatness-filtering
algorithm. \Cref{sec:method-coinflip} describes the evaluation
protocol and the rational-attacker floor applied to reported ASR.
\Cref{sec:method-protocol} lists the locked parameters used by every
experiment in \Cref{ch:results}.

\section{Overview}
\label{sec:method-overview}

For each real password $w^{*}$, a generator produces $k-1$ candidate
honeywords which are combined with $w^{*}$ into a sweetword set
$S = \{w_1, \ldots, w_k\}$ of size $k = 20$. Two independently trained
character-level CNNs interact with $S$ at different points in the
pipeline. The first, the \emph{oracle}, is used at storage time by the
filtering algorithm to decide whether $S$ satisfies the flatness
conditions defined in \Cref{sec:method-formalism} and, if not, to
trigger regeneration. The second, the \emph{attacker}, is used at
evaluation time to measure the attacker success rate on the accepted
sets. The two networks share the same architecture but are trained on
disjoint corpora with different random seeds
(\Cref{sec:method-cnn}), which is what allows us to report ASR against
an attacker that is not the model the filter was tuned against.

Filtering is performed once, at password storage time, and is never
invoked during authentication. The construction is generator-agnostic:
we apply the same filter, with the same thresholds, on top of ten
distinct honeyword generators, including one (HoneyGen Baseline) that
performs no filtering and serves as the reference point for measuring
the effect of the filter.

\section{Attacker-Relative Set-Level Flatness}
\label{sec:method-formalism}

\subsection{Sweetword Set and Attacker Model}

A sweetword set is a size-$k$ multiset
$S = \{w_1, \ldots, w_k\}$ in which exactly one element $w^{*}$ is the
real password and the remaining $k-1$ elements are honeywords produced
by a generator. The attacker $A$ is a binary classifier trained on
labeled examples of the form $(w, y)$ with $y = 1$ for real passwords
and $y = 0$ for honeywords. For every candidate $w_i$, the classifier
outputs an independent score
%
\begin{equation}
s_i \;=\; s(w_i) \;=\; P_A(y = 1 \mid w_i) \in [0, 1].
\end{equation}
%
The scores are per-candidate probabilities from a binary classifier,
not a distribution over $S$. In particular
$\sum_{i} s_i \neq 1$ in general, and the classifier processes each
$w_i$ without seeing any other member of $S$.

\subsection{Set-Level Induced Distribution}

The attacker's actual decision problem is comparative: given all of
$S$, choose one candidate. We model this comparison by renormalizing
the per-candidate scores to form a distribution over $S$:
%
\begin{equation}
q_i \;=\; \frac{s_i}{\sum_{j=1}^{k} s_j}, \qquad
\sum_{i=1}^{k} q_i \;=\; 1.
\label{eq:q}
\end{equation}
%
The vector $q = (q_1, \ldots, q_k)$ is the \emph{attacker-induced
distribution} on $S$. Under ideal flatness $q$ is uniform, so
$q_i = 1/k$ for every $i$. We write $q_{\text{real}}$ for the entry of
$q$ corresponding to $w^{*}$ and use it as the central quantity that
the filter constrains.

\subsection{Flatness Metrics}

Four quantities computed from the pair $(s, q)$ characterize how far
$S$ deviates from flatness under $A$.

\paragraph{Rank-based flatness.}
Order the candidates by $s_i$ in decreasing order and let
$\text{rank}(w^{*})$ be the position of $w^{*}$ in that order. We
define
%
\begin{equation}
F_{\text{rank}} \;=\; \frac{\text{rank}(w^{*}) - 1}{k - 1} \in [0, 1].
\end{equation}
%
A value of $0$ means the attacker ranks $w^{*}$ first (worst case for
the defender); a value of $1$ means it ranks $w^{*}$ last; the
expected value under a uniform ranker is $1/2$.

\paragraph{Normalized entropy.}
The Shannon entropy of $q$ divided by its maximum $\log k$:
%
\begin{equation}
F_{\text{ent}} \;=\; \frac{H(q)}{\log k}
\;=\; -\frac{1}{\log k}\sum_{i=1}^{k} q_i \log q_i.
\end{equation}
%
$F_{\text{ent}} = 1$ if and only if $q$ is uniform. Entropy captures
global smoothness of $q$ but does not by itself constrain which entry
of $q$ is the largest, and a set can have $F_{\text{ent}}$ close to
$1$ while still assigning the top rank to $w^{*}$. We therefore treat
$F_{\text{ent}}$ as a secondary stabilizing condition rather than a
primary one.

\paragraph{Gap metrics.}
Let $s_{(1)} \ge s_{(2)} \ge \cdots \ge s_{(k)}$ be the sorted raw
scores. The top gap
%
\begin{equation}
\Delta_{\text{top}} \;=\; s_{(1)} - s_{(2)}
\end{equation}
%
measures how strongly the attacker prefers its top candidate over the
runner-up, and the real gap
%
\begin{equation}
\Delta_{\text{real}} \;=\; s_{(1)} - s(w^{*})
\end{equation}
%
measures how far the real password falls below the top score. Gap
metrics are computed on raw scores rather than on $q$ because ranking
is invariant under the positive rescaling of \Cref{eq:q} and the raw
scores preserve the absolute magnitude of the classifier's preference.

\paragraph{$\varepsilon$-flatness.}
Let $q_{\max} = \max_i q_i$ denote the largest entry of the induced
distribution. The set $S$ is \emph{$\varepsilon$-flat} under $A$ if
%
\begin{equation}
q_{\max} \;\le\; \frac{1}{k} + \varepsilon.
\label{eq:eps-flat}
\end{equation}
%
Setting $\varepsilon = 0$ forces $q_{\max} \le 1/k$, which together
with $\sum_i q_i = 1$ recovers exact uniformity of the whole
distribution; increasing $\varepsilon$ relaxes the bound. Because
$q_{\text{real}} \le q_{\max}$, this constraint upper-bounds the
attacker's posterior on every sweetword, including the true credential,
and is the primary security constraint enforced by the filter. This
additive slack $\varepsilon$ on the maximum entry differs from the
classical multiplicative bound $(1+\varepsilon)/k$ of
\Cref{sec:bg-flatness}: the two coincide at $\varepsilon = 0$, but for
$\varepsilon > 0$ the additive form places the threshold at
$1/k + \varepsilon$ rather than $(1+\varepsilon)/k$, so the same
numerical $\varepsilon$ is a looser bound here. We adopt the additive
form because it acts directly on the peak of the attacker-induced
distribution over the set, whereas the classical bound is stated on the
attacker's first-guess success probability.

\paragraph{Attacker success rate.}
Given the ranking induced by $s$, the attacker success rate at rank
$t$ is
%
\begin{equation}
\mathrm{ASR}@t \;=\; \Pr\!\left[\, \text{rank}(w^{*}) \le t \,\right],
\end{equation}
%
where the probability is taken over the sample of evaluation sets.
Under ideal flatness $\mathrm{ASR}@t = t/k$. We report
$\mathrm{ASR}@1$, $\mathrm{ASR}@3$, and $\mathrm{ASR}@5$ throughout
\Cref{ch:results}, with $\mathrm{ASR}@5$ as the primary security
metric.

\subsection{Operational Definition of Flatness}

The set $S$ is declared flat under the attacker $A$ if it satisfies
\Cref{eq:eps-flat} together with two auxiliary conditions on
$F_{\text{ent}}$ and $\Delta_{\text{top}}$. The three thresholds and
their justification are given in \Cref{sec:method-protocol}. Sets that
fail any of the three conditions are rejected by the filter of
\Cref{sec:method-filter}.

\section{Corpora and Splits}
\label{sec:method-corpora}

\subsection{Sources}

The corpus is the cleaned RockYou leak, from which we draw disjoint
oracle and attacker subsets after deduplication, encoding validation,
and truncation to the CNN's input length. The oracle corpora are the
three training sizes $N_{\text{train}} \in \{2000,\,5000,\,10000\}$
used throughout \Cref{ch:results}; the attacker corpus is a fixed
disjoint subset of $5000$ passwords, split $80/10/10$ and reused
unchanged for every oracle size (the training composition is detailed
in \Cref{sec:method-cnn}). All entries are unprocessed leaked passwords; no case folding,
character substitution, or semantic normalization is applied beyond
the length truncation. Extending the benchmark to a second corpus such
as \emph{darkc0de} is left as future work (\Cref{ch:conclusion}).

\subsection{Disjoint Oracle and Attacker Corpora}

The pipeline uses two disjoint corpora. The \emph{oracle corpus}
supplies the real passwords $w^{*}$ from which sweetword sets are
built, and also provides the training data for the oracle CNN used by
the filter; a separate oracle is trained for each generator size
$N_{\text{train}}$. The \emph{attacker corpus}, drawn from the same
source but disjoint from the oracle corpus, provides the training data
for the attacker CNN used to measure $\mathrm{ASR}$ on the accepted
sets. Unlike the oracle, the attacker is trained once on the training
split of a fixed $5000$-password attacker corpus (\Cref{sec:method-cnn})
and reused across every generator size, so that a single, consistent
adversary evaluates all of them. Because
the attacker is trained on passwords that neither the generators nor
the oracle have seen, its evaluation performance reflects
generalization to unseen inputs rather than memorization of the sets it
will later be asked to rank. This choice raises the baseline
$\mathrm{ASR}$ compared to a shared-corpus setup, but the resulting
numbers are the ones that matter for a real deployment.

\subsection{Set-Level Splits}

Within the oracle corpus, real passwords are partitioned into
training, validation, and test subsets. Sweetword sets are constructed
for every real password in each subset, and all sets are concatenated
into a single evaluation stream with globally unique set identifiers.
This \texttt{EVAL\_FULL = True} configuration is required to avoid
staircase artifacts in the ASR curve at small evaluation sizes, which
would otherwise arise from evaluating only the tail of a single split.

\section{The CNN Discriminator}
\label{sec:method-cnn}

\subsection{Architecture}

The classifier used for both oracle and attacker roles is a
character-level convolutional network. Each password is truncated or
right-padded to a fixed length and mapped through a character
embedding of dimension $64$. Two one-dimensional convolutional layers
with ReLU activations extract local $n$-gram features, followed by
max-pooling to reduce dimensionality. A fully connected layer projects
to a two-dimensional output over which a softmax produces
$P(\text{real} \mid w)$ and $P(\text{honey} \mid w)$. The architecture
is inspired by the learned-attacker paradigm of
\textcite{dani2024passfilter} but is trained independently and does
not reuse any released weights.

\subsection{Training}

The classifier is trained as a binary classifier with cross-entropy
loss
%
\begin{equation}
\mathcal{L} \;=\; -\bigl(y \log s \;+\; (1 - y) \log (1 - s)\bigr).
\end{equation}
%
Positive examples are real passwords and negative examples are
honeywords produced by the HoneyGen Baseline generator described in
\Cref{subsec:method-honeygen}, which combines the FastText password
model with rule-based tweaking. Training the classifier against
HoneyGen negatives calibrates it to the dominant family of decoys
studied in this thesis and produces a strong and realistic adversary.
The oracle is trained on the training split of its own
$N_{\text{train}}$ corpus. The attacker is trained once on a fixed
$5000$-password subset of RockYou, disjoint from every oracle corpus
and split $80/10/10$ into training, validation, and test partitions;
its training set is built by running HoneyGen Baseline with $k = 20$ on
the $4000$ real passwords of the training partition, producing $19$
honeywords per real password. The attacker CNN is therefore trained on
$4000$ real and $76{,}000$ honeyword examples, and this single trained
attacker is reused to evaluate every generator at every training size.
Training uses batch size $\text{BS\_CNN} = 256$ and runs for ten epochs
for both oracle and attacker roles.

\subsection{Dual Role and Seed Separation}

The two networks share the architecture and training procedure just
described but differ in three respects. First, they are trained on
disjoint corpora (\Cref{sec:method-corpora}). Second, the
attacker uses a seed offset of $7777$ relative to the oracle, so the
two networks are not identical even where their training procedures
coincide. Third, only the oracle interacts with the filter at storage
time; the attacker is invoked exclusively at evaluation time and never
observes the acceptance decisions the oracle made.

\subsection{Robustness Attackers}
\label{sec:method-robustness-attackers}

The CNN above is the attacker the filter is tuned against. To test
whether the flatness it enforces survives a change of attacker, we
evaluate the accepted sets against two further attackers that the filter
never saw. Both are used only at evaluation time, on sets the filter has
already accepted; neither takes part in filtering.

\paragraph{Frequency ranker.}
The frequency attacker is a classical, non-learned baseline. It first
builds a frequency table from the attacker corpus by counting how often
each password occurs. To rank a sweetword set it scores each of the $k$
candidates by that count --- a common password scores high, a rare or
never-seen one scores low (unseen candidates receive the lowest score)
--- and then guesses the candidates from most to least frequent, on the
assumption that a real password is more likely to be a popular one. It
uses no neural network and no training beyond counting. This is exactly
why it is a useful control: it is the ranker-style attacker that most
earlier honeyword evaluations assumed, so comparing it against the CNN
shows how much a learned discriminator adds over simple popularity.

\paragraph{Mixed-corpus CNN.}
The mixed-corpus attacker has the same architecture and training
procedure as the oracle CNN, but its negative examples are drawn from
several generator families together --- Chaffing-by-Tweaking, the
Password-Model, HoneyGen, and the GAN --- instead of from HoneyGen
Baseline alone. Training on a broader mix stops the attacker from
over-fitting to one decoy family and lets it flag sets, such as the
GAN's, that a single-family CNN misses (\Cref{sec:results-gan}).

\section{Honeyword Generation Methods}
\label{sec:method-gens}

We benchmark ten honeyword generators. Six of them are unfiltered
baselines that expose how each generation strategy behaves against the
learned attacker; four are filtered variants that combine a generator
with the pre-screening or rejection sampling stages of
\Cref{sec:method-filter}.

\subsection{Chaffing Family}
\label{subsec:method-chaffing}

The three \emph{chaffing} methods derive candidates directly from
$w^{*}$ and do not perform any oracle-gated selection.

\paragraph{Chaffing-by-Tweaking.}
Rule-based edits applied to $w^{*}$: character-class substitutions on
letters, digits, and symbols; suffix and prefix modifications; and
short case flips. Candidates preserve the length structure of the real
password.

\paragraph{Chaffing-with-Password-Model.}
Nearest neighbors of $w^{*}$ in the FastText embedding space, using
the skipgram model with \texttt{minCount} = 1 and \texttt{minn} = 2
trained on the oracle corpus. Candidates share sub-word structure with
$w^{*}$ but need not preserve its length.

\paragraph{Chaffing-with-Hybrid-Model.}
This generator combines the two streams above. For a real password
$w^{*}$ it draws candidates from both sources -- the length-preserving
rule-based edits of Chaffing-by-Tweaking and the FastText nearest
neighbors of Chaffing-with-Password-Model -- and forms the $k-1$
decoys of each sweetword set from this combined pool. A resulting set
can therefore contain both surface-level variants of $w^{*}$ (from the
tweaking stream, which keep its length and character structure) and
sub-word-similar neighbors (from the password-model stream, which share
its morphology but not necessarily its length). Mixing the two makes
the accepted sets harder for an attacker to characterize with a single
heuristic, since no one edit rule or embedding-distance cutoff explains
all the decoys. Because both streams reuse the same underlying
components as HoneyGen Baseline, the random seed that governs how the
two are combined is offset by $5000$ from the base seed; this ensures
the hybrid method produces a distinct sample rather than reproducing
the sets of HoneyGen Baseline.

\subsection{HoneyGen Family}
\label{subsec:method-honeygen}

The five \emph{HoneyGen} methods build on the algorithm of
\textcite{dionysiou2021honeygen}, which combines the FastText password
model with rule-based tweaking and applies frequency-based pruning to
the resulting candidate pool.

\paragraph{HoneyGen Baseline.}
The HoneyGen algorithm as described by
\textcite{dionysiou2021honeygen}, without any oracle interaction. This
is the reference point against which the effect of every filter is
measured.

\paragraph{HoneyGen Set-Filter.}
HoneyGen Baseline followed by set-level rejection sampling only.
Candidate honeywords are used as produced; the oracle decides
acceptance at the sweetword-set level.

\paragraph{HoneyGen Prescreen-Add and Prescreen-Ratio.}
HoneyGen Baseline restricted to its password-model component (the
\texttt{honeygen-model-only} configuration), followed by both
candidate-level pre-screening and set-level rejection sampling. The
Add variant uses an additive band condition
$|s(w) - s(w^{*})| \le \tau_{\text{band}}$, and the Ratio variant uses
a multiplicative band condition
$s(w) / s(w^{*}) \le \tau_{\text{band-ratio}}$.

\paragraph{HoneyGen Hybrid-Add and Hybrid-Ratio.}
Identical to the Prescreen variants except that the candidate stream
is the full HoneyGen hybrid (password model plus tweaking) rather than
password-model only. Both use the two-stage filter.

\subsection{GAN}

A generator adversarial network trained separately for each
evaluation-set size. Candidates are drawn from the trained generator,
followed by oversampling by a factor of approximately five,
deduplication, and removal of malformed entries; a random subset of
the survivors is retained as the $k-1$ honeywords. The GAN is not
gated by the oracle. Its results in \Cref{ch:results} require a
qualification that we discuss there.

\section{Hybrid Flatness-Filtering Algorithm}
\label{sec:method-filter}

The filter has two stages. Candidate-level pre-screening removes
individually distinguishable decoys before a sweetword set is
assembled, and set-level rejection sampling discards sets whose
induced distribution violates one of the flatness conditions from
\Cref{sec:method-formalism}. Both stages use the oracle CNN of
\Cref{sec:method-cnn}.

\subsection{Candidate-Level Pre-Screening}

For each candidate honeyword $w$, the oracle produces a score
$s(w) \in [0, 1]$. The pre-screening condition compares this score to
the score $s(w^{*})$ of the real password within a band around it:
the \emph{additive} band accepts $w$ when
%
\begin{equation}
|\, s(w) - s(w^{*}) \,| \;\le\; \tau_{\text{band}},
\end{equation}
%
and the \emph{multiplicative} band accepts $w$ when its score lies
within a factor $\tau_{\text{band-ratio}}$ of $s(w^{*})$ on either side,
%
\begin{equation}
\frac{1}{\tau_{\text{band-ratio}}}\, s(w^{*}) \;\le\; s(w) \;\le\;
\tau_{\text{band-ratio}}\, s(w^{*}),
\qquad \tau_{\text{band-ratio}} > 1.
\end{equation}
%
The two variants correspond to the Prescreen and Hybrid families of
\Cref{sec:method-gens}. Both bands are two-sided: they remove
candidates that the oracle is confident about in either direction --
whether obviously honeyword-looking, with a score far below
$s(w^{*})$, or obviously real-looking, with a score far above it --
since either kind stands out in the induced distribution. One practical
adjustment is needed for the multiplicative band: when the oracle is
undertrained it sometimes scores the real password very close to zero,
which would shrink the band to almost nothing and reject every
candidate. To avoid this, the band is built around
$\max(s(w^{*}), 0.05)$ instead of $s(w^{*})$ directly, so that a
near-zero real score does not collapse the band.

\subsection{Set-Level Rejection Sampling}

After $k-1$ candidates survive pre-screening, the sweetword set
$S = \{w^{*}\} \cup \{w_1, \ldots, w_{k-1}\}$ is passed to the oracle,
which produces the score vector $s = (s_1, \ldots, s_k)$ and, via
\Cref{eq:q}, the induced distribution $q$. The set is accepted if all
three of the following hold:
%
\begin{align}
q_{\max} &\;\le\; \frac{1}{k} + \varepsilon_{\max}, \\
F_{\text{ent}}  &\;\ge\; 1 - \delta_{H}, \\
\Delta_{\text{top}} &\;\le\; \text{gap}_{\max}.
\end{align}
%
If any condition fails, the entire set is discarded and a fresh one is
generated. Regeneration continues until acceptance or until a
regeneration cap $R_{\max}$ is reached, in which case the best set
observed so far (measured by minimum $q_{\max}$) is retained.

There is one design limitation of this test worth stating explicitly.
The predicate caps only the largest entry $q_{\max}$; it imposes no
lower bound on any entry. A set in which some honeyword receives an
anomalously low $q$ can therefore still pass, even though the resulting
distribution is not perfectly uniform. In practice this asymmetry does
not affect ASR, because ASR responds to the top of the ranking rather
than the tail; we return to the point in \Cref{ch:results}.

\subsection{Algorithm}

The combined procedure is stated in \Cref{alg:hff}. The pseudocode
uses generic band operator $\text{band}(\cdot)$ that stands for either
the additive or the ratio variant of the pre-screening rule.

\begin{algorithm}[H]
\caption{Hybrid Flatness Filtering}
\label{alg:hff}
\begin{algorithmic}[1]
\Require real password $w^{*}$, generator $G$, oracle $A$,
    thresholds $\tau_{\text{band}},\ \varepsilon_{\max},\
    \delta_{H},\ \text{gap}_{\max}$, cap $R_{\max}$
\Ensure sweetword set $S$
\State $r \gets 0$;\ \ $S^{\text{best}} \gets \varnothing$;\ \ 
    $q^{\text{best}} \gets +\infty$
\While{$r < R_{\max}$}
    \State $C \gets G.\textsc{Generate}(w^{*},\ N_{\text{over}})$
    \Comment{oversampled candidate pool}
    \State $C \gets \{\, w \in C : \textsc{Band}(s(w), s(w^{*})) \,\}$
    \Comment{candidate pre-screening}
    \If{$|C| < k - 1$}
        \State $r \gets r + 1$;\ \textbf{continue}
    \EndIf
    \State $S \gets \{w^{*}\} \cup \textsc{TakeFirst}(C,\ k-1)$
    \State compute $s_i = A(w_i)$ and $q_i$ via \Cref{eq:q}
    \State compute $q_{\max}$, $F_{\text{ent}}$, and $\Delta_{\text{top}}$
    \If{$q_{\max} \le 1/k + \varepsilon_{\max}$ \textbf{and}
        $F_{\text{ent}} \ge 1 - \delta_{H}$ \textbf{and}
        $\Delta_{\text{top}} \le \text{gap}_{\max}$}
        \State \Return $S$
    \EndIf
    \If{$q_{\max} < q^{\text{best}}$}
        \State $S^{\text{best}} \gets S$;\ 
            $q^{\text{best}} \gets q_{\max}$
    \EndIf
    \State $r \gets r + 1$
\EndWhile
\State \Return $S^{\text{best}}$
\end{algorithmic}
\end{algorithm}

\subsection{Complexity}

Let $C_{\text{CNN}}$ denote the cost of one forward pass through the
oracle and let $\mathbb{E}[R]$ denote the expected number of
regeneration attempts per password. The expected cost of producing one
accepted sweetword set is
%
\begin{equation}
\mathcal{O}\!\left(\, \mathbb{E}[R] \cdot k \cdot C_{\text{CNN}} \,\right),
\end{equation}
%
plus the cost of the generator itself. Since $C_{\text{CNN}}$
dominates the arithmetic cost of any of the generators used here, the
filter's total cost is close to $\mathbb{E}[R] \cdot k$ oracle forward
passes per password. Batched scoring keeps this modest for
$k = 20$ and the $R_{\max}$ values reported in
\Cref{sec:method-protocol}, and filtering happens once at storage
time rather than at every authentication.

\section{Evaluation Protocol}
\label{sec:method-coinflip}

\subsection{Design}

Each generator is trained once at a fixed training size
$N_{\text{train}}$, and every evaluation set is produced by that single
trained generator. $\mathrm{ASR}@t$ is then computed over a test set of
sweetword sets, restricted to the sets the filter accepted
(\Cref{sec:method-filter}). Because the generator is fixed for a given
$N_{\text{train}}$, the reported figures compare methods on the same
footing rather than conflating a change of generator with a change of
measurement. We report results at three training sizes,
$N_{\text{train}} \in \{2000,\,5000,\,10000\}$, evaluating each on a
test set of $2000$ sweetword sets.

\subsection{Rational Attacker Floor}

The CNN attacker's raw $\mathrm{ASR}@t$ --- the success of following the
CNN's $t$ highest-scored candidates --- can fall \emph{below} the
random-guessing baseline $t/k$ on some sets. This is not a sign that the
set is safe. It means the CNN's ranking is \emph{anti-informative}: it
places the real password near the \emph{bottom} of the order, so
following its top guesses does worse than chance. A rational attacker
who observed this would not keep following the CNN; it would invert its
strategy and guess the candidates the CNN scores \emph{lowest}, which
then succeeds far more often than chance.

We therefore report the success of a rational attacker that uses the CNN
ranking in whichever direction helps, falling back on random guessing
when neither direction is informative:
%
\begin{equation}
\mathrm{ASR}_{\text{reported}}(t)
    \;=\; \max\!\bigl(\, \mathrm{ASR}_{\text{CNN}}(t),\
    \mathrm{ASR}_{\overline{\text{CNN}}}(t),\ t/k \,\bigr),
\label{eq:rational}
\end{equation}
%
where $\mathrm{ASR}_{\text{CNN}}(t)$ counts the sets whose real password
is among the CNN's top $t$ candidates,
$\mathrm{ASR}_{\overline{\text{CNN}}}(t)$ counts those whose real
password is among the CNN's \emph{bottom} $t$ candidates (the inverted
strategy), and $t/k$ is plain random guessing. A genuinely flat
generator drives all three terms to $t/k$, so the reported value is
$t/k$. A generator whose real password the CNN systematically ranks low
--- notably the GAN of \Cref{sec:results-gan} --- is caught by the
inverted term, which reports the high success such a set really allows
instead of masking it at $t/k$. This is the value used in every ASR
figure and table in \Cref{ch:results}.

\section{Experimental Protocol and Parameters}
\label{sec:method-protocol}

\subsection{Locked Parameters}

The following parameters are fixed for every experiment reported in
\Cref{ch:results}. The flatness thresholds were set in two stages. An
initial, strict configuration ($\varepsilon_{\max} = 0.02$,
$\delta_{H} = 0.04$, $\text{gap}_{\max} = 0.01$) was chosen from the
theoretical rationale in \Cref{sec:method-formalism} and preliminary
runs on the training split. As reported in
\Cref{sec:method-threshold-tradeoff}, this configuration drove the
ratio-based filters to acceptance rates below $10\%$, at which point
the filtered methods can no longer be meaningfully compared to the
unfiltered baselines. We therefore adopted a relaxed configuration,
listed below, that restores usable acceptance across all filters at
every corpus size. None of the thresholds was tuned on the evaluation
split; the revision was driven entirely by the acceptance-rate
diagnostic on the training split.

\begin{center}
\begin{tabular}{lll}
\toprule
\textbf{Symbol} & \textbf{Value} & \textbf{Role} \\
\midrule
$k$                       & $20$    & Sweetword set size \\
$\varepsilon_{\max}$      & $0.05$  & $\varepsilon$-flatness bound (see note below) \\
$\delta_{H}$              & $0.10$  & Entropy slack, $F_{\text{ent}} \ge 1 - \delta_{H}$ \\
$\text{gap}_{\max}$       & $0.05$  & Top-gap threshold \\
$\tau_{\text{band}}$      & $0.15$  & Additive pre-screening band \\
$\tau_{\text{band-ratio}}$& $2.0$   & Multiplicative pre-screening band \\
$R_{\max}$                & $50$    & Regeneration cap \\
$\text{BS\_CNN}$          & $256$   & CNN training batch size \\
Epochs (oracle)           & $10$    & CNN oracle training epochs \\
Epochs (attacker)         & $10$    & CNN attacker training epochs \\
$\text{SEED}$             & $42$    & Base random seed \\
Attacker seed offset      & $7777$  & Added to base seed for attacker \\
Hybrid model seed offset  & $5000$  & Added to base seed for Hybrid-Model \\
\bottomrule
\end{tabular}
\end{center}

A single value of $\varepsilon_{\max} = 0.05$ is used throughout, both
in the definition of $\varepsilon$-flatness (\Cref{eq:eps-flat}) and in
the filter's acceptance predicate (\Cref{sec:method-filter}), which
share the additive form $q_{\max} \le 1/k + \varepsilon_{\max}$.

\subsection{Threshold Selection and the Accept-Rate / ASR Trade-off}
\label{sec:method-threshold-tradeoff}

The flatness thresholds control a trade-off between two quantities that
a deployed countermeasure must satisfy simultaneously. A stricter gate
produces flatter accepted sets, lowering $\mathrm{ASR}$, but rejects a
larger share of passwords; each rejected password reaches the
regeneration cap without an accepted set and is excluded from the
reported $\mathrm{ASR}$, which is therefore conditional on acceptance.
The accept rate measures how often the filter succeeds, and the
$\mathrm{ASR}$ measures how flat the sets are when it does.

Under the initial thresholds, the ratio-based filters accepted fewer
than $10\%$ of candidate passwords at $N_{\text{train}} \in \{5000,
10000\}$ (Set-Filter reached $1.2\%$ at $N_{\text{train}} = 5000$),
which makes their $\mathrm{ASR}$ a statement about a small and
self-selected minority of passwords rather than about the method as a
deployed defense. The adopted thresholds raise every filter above
$25\%$ acceptance at every corpus size (\Cref{ch:results}), at the cost
of a modest increase in $\mathrm{ASR}$ on the accepted sets. This is a
deliberate choice: a filter that protects a single-digit percentage of
users is not a usable countermeasure regardless of how flat its
accepted sets are, so acceptance is treated as a first-class
deployability requirement rather than a secondary diagnostic.

Consequently, methods are compared in two stages in \Cref{ch:results}.
A method must first clear a usable acceptance threshold; only among the
methods that do so is the flattest (lowest-$\mathrm{ASR}$) one preferred.
Under this criterion the ratio-based filters are the recommended
configurations, since they deliver the flattest sets among the methods
that remain deployable, whereas the additive filters accept more
passwords but sit closer to the unfiltered baseline in flatness.

\subsection{Metrics and Statistical Treatment}

For each generator and each $N$ we report $\mathrm{ASR}@1$,
$\mathrm{ASR}@3$, and $\mathrm{ASR}@5$ using
\Cref{eq:rational}, with $\mathrm{ASR}@5$ as the primary metric.
Alongside each ASR value we report the mean and median of
$F_{\text{rank}}$, $F_{\text{ent}}$, $q_{\text{real}}$, and
$\Delta_{\text{top}}$ over the evaluation sets, together with the
number of accepted sets for the filtered generators. Confidence
intervals at small $N$ are wide, approximately $\pm 20$ percentage
points at $N \le 500$, and we do not draw conclusions from small-$N$
comparisons that fall inside this margin.

\subsection{Robustness Check}

The primary corpus for all experiments is RockYou, and the primary
attacker is the HoneyGen-tuned CNN. Robustness is assessed by
re-evaluating the accepted sets against two further attacker classes
--- a mixed-corpus CNN trained on several decoy families and a
frequency ranker --- as reported in \Cref{sec:results-freq}. Extending
the benchmark to a second corpus such as \emph{darkc0de} is left as
future work (\Cref{ch:conclusion}).
